\documentclass[11pt,a4paper]{article}
% preamble.tex — estilo común de todos los apuntes Froth.
% Cada apunte hace % preamble.tex — estilo común de todos los apuntes Froth.
% Cada apunte hace % preamble.tex — estilo común de todos los apuntes Froth.
% Cada apunte hace \input{preamble} justo después de \documentclass.
\usepackage[margin=2.4cm]{geometry}
\usepackage{xcolor}
\definecolor{litblue}{RGB}{31,79,121}
\definecolor{litgreen}{RGB}{27,94,32}
\definecolor{litorange}{RGB}{230,81,0}
\usepackage{parskip}
\usepackage{titlesec}
\titleformat{\section}{\Large\bfseries\color{litblue}}{\thesection}{0.6em}{}
\titleformat{\subsection}{\large\bfseries\color{litblue}}{\thesubsection}{0.6em}{}
\usepackage{tcolorbox}
\usepackage{fancyhdr}
\pagestyle{fancy}
\fancyhf{}
\fancyhead[L]{\small\color{litblue}Froth — Apuntes de ML}
\fancyhead[R]{\small\thepage}
\renewcommand{\headrulewidth}{0.4pt}
\usepackage[colorlinks=true,linkcolor=litblue,urlcolor=litblue]{hyperref}

% Cabecera de cada apunte: \cabecera{numero}{titulo}
\newcommand{\cabecera}[2]{%
  \begin{tcolorbox}[colback=litblue,coltext=white,colframe=litblue,arc=2mm]
    {\Large\bfseries Apunte #1 — #2}\\[3pt]
    {\small Proyecto Froth · aprender Machine Learning construyendo}
  \end{tcolorbox}\vspace{0.8em}}

% Cajas temáticas
\newtcolorbox{concepto}[1]{colback=blue!4,colframe=litblue,title={Concepto: #1},arc=1.5mm}
\newtcolorbox{practica}{colback=green!5,colframe=litgreen,title={Pruébalo tú (teclado en mano)},arc=1.5mm}
\newtcolorbox{ojo}{colback=orange!8,colframe=litorange,title={Ojo / error típico},arc=1.5mm}
\newtcolorbox{resumen}{colback=gray!8,colframe=gray!60,title={En una frase},arc=1.5mm}

\newcommand{\codigo}[1]{\texttt{#1}}
 justo después de \documentclass.
\usepackage[margin=2.4cm]{geometry}
\usepackage{xcolor}
\definecolor{litblue}{RGB}{31,79,121}
\definecolor{litgreen}{RGB}{27,94,32}
\definecolor{litorange}{RGB}{230,81,0}
\usepackage{parskip}
\usepackage{titlesec}
\titleformat{\section}{\Large\bfseries\color{litblue}}{\thesection}{0.6em}{}
\titleformat{\subsection}{\large\bfseries\color{litblue}}{\thesubsection}{0.6em}{}
\usepackage{tcolorbox}
\usepackage{fancyhdr}
\pagestyle{fancy}
\fancyhf{}
\fancyhead[L]{\small\color{litblue}Froth — Apuntes de ML}
\fancyhead[R]{\small\thepage}
\renewcommand{\headrulewidth}{0.4pt}
\usepackage[colorlinks=true,linkcolor=litblue,urlcolor=litblue]{hyperref}

% Cabecera de cada apunte: \cabecera{numero}{titulo}
\newcommand{\cabecera}[2]{%
  \begin{tcolorbox}[colback=litblue,coltext=white,colframe=litblue,arc=2mm]
    {\Large\bfseries Apunte #1 — #2}\\[3pt]
    {\small Proyecto Froth · aprender Machine Learning construyendo}
  \end{tcolorbox}\vspace{0.8em}}

% Cajas temáticas
\newtcolorbox{concepto}[1]{colback=blue!4,colframe=litblue,title={Concepto: #1},arc=1.5mm}
\newtcolorbox{practica}{colback=green!5,colframe=litgreen,title={Pruébalo tú (teclado en mano)},arc=1.5mm}
\newtcolorbox{ojo}{colback=orange!8,colframe=litorange,title={Ojo / error típico},arc=1.5mm}
\newtcolorbox{resumen}{colback=gray!8,colframe=gray!60,title={En una frase},arc=1.5mm}

\newcommand{\codigo}[1]{\texttt{#1}}
 justo después de \documentclass.
\usepackage[margin=2.4cm]{geometry}
\usepackage{xcolor}
\definecolor{litblue}{RGB}{31,79,121}
\definecolor{litgreen}{RGB}{27,94,32}
\definecolor{litorange}{RGB}{230,81,0}
\usepackage{parskip}
\usepackage{titlesec}
\titleformat{\section}{\Large\bfseries\color{litblue}}{\thesection}{0.6em}{}
\titleformat{\subsection}{\large\bfseries\color{litblue}}{\thesubsection}{0.6em}{}
\usepackage{tcolorbox}
\usepackage{fancyhdr}
\pagestyle{fancy}
\fancyhf{}
\fancyhead[L]{\small\color{litblue}Froth — Apuntes de ML}
\fancyhead[R]{\small\thepage}
\renewcommand{\headrulewidth}{0.4pt}
\usepackage[colorlinks=true,linkcolor=litblue,urlcolor=litblue]{hyperref}

% Cabecera de cada apunte: \cabecera{numero}{titulo}
\newcommand{\cabecera}[2]{%
  \begin{tcolorbox}[colback=litblue,coltext=white,colframe=litblue,arc=2mm]
    {\Large\bfseries Apunte #1 — #2}\\[3pt]
    {\small Proyecto Froth · aprender Machine Learning construyendo}
  \end{tcolorbox}\vspace{0.8em}}

% Cajas temáticas
\newtcolorbox{concepto}[1]{colback=blue!4,colframe=litblue,title={Concepto: #1},arc=1.5mm}
\newtcolorbox{practica}{colback=green!5,colframe=litgreen,title={Pruébalo tú (teclado en mano)},arc=1.5mm}
\newtcolorbox{ojo}{colback=orange!8,colframe=litorange,title={Ojo / error típico},arc=1.5mm}
\newtcolorbox{resumen}{colback=gray!8,colframe=gray!60,title={En una frase},arc=1.5mm}

\newcommand{\codigo}[1]{\texttt{#1}}

\begin{document}
\cabecera{17}{El filtro de relevancia: el sistema se limpia solo}

\section{El problema que TÚ detectaste}

Validando el mapa (paso 3.7) viste un cluster etiquetado \emph{``women, men, commodity,
masculinities''}. Diagnóstico: los papers infiltrados de la Fase 1 (poesía, telecomunicaciones,
historia --- colisiones de siglas en las búsquedas) dominaban las keywords de un cluster que
también contenía papers legítimos de economía del litio. La causa raíz estaba en el
\emph{pull}, no en el mapa.

\section{La solución elegante: medir cada paper contra el TOPIC}

\begin{concepto}{score de relevancia}
Ya tenemos todas las piezas: cada paper tiene su vector, y el \textbf{título de tu tesis}
también puede tenerlo. Entonces: \emph{relevancia de un paper} $=$ similitud coseno entre su
vector y el vector del TOPIC. Un paper de poesía queda lejos del título de una tesis de
flotación $\rightarrow$ score bajo $\rightarrow$ fuera. El sistema usa \textbf{sus propios
embeddings para limpiarse a sí mismo} --- cero reglas manuales, cero listas negras.
\end{concepto}

\begin{concepto}{el umbral lo dictan los datos, no la intuición}
¿Dónde cortar? Miramos la distribución real de los 126 scores:
\begin{center}
\begin{tabular}{ll}
Infiltrados (poesía, telecom, historia...) & 0.495 -- 0.543 \\
\textbf{--- hueco natural ---} & \\
Todo lo legítimo (incluida economía del Li) & 0.587 -- 0.90 \\
\end{tabular}
\end{center}
Hay un \textbf{hueco} entre 0.543 y 0.587: el umbral \codigo{RELEVANCE\_THRESHOLD = 0.55}
cae en medio y separa perfecto. Así se eligen umbrales en serio: mirando dónde los datos ya
están separados, no inventando un número bonito.
\end{concepto}

\section{Dónde vive en el código}

\begin{itemize}
  \item \codigo{embed.py} $\rightarrow$ \codigo{relevance\_to\_topic()}: calcula los scores y
        los guarda como columna \codigo{relevance} en \codigo{papers.parquet}.
  \item \codigo{config.py} $\rightarrow$ \codigo{RELEVANCE\_THRESHOLD = 0.55} (con el
        razonamiento comentado).
  \item \codigo{cluster.py} $\rightarrow$ filtra ANTES de UMAP: la basura ni siquiera
        participa en el acomodo, así el mapa entero mejora.
\end{itemize}

\section{Qué cambió en el mapa (iteración real de data science)}

\begin{itemize}
  \item Las etiquetas vergonzosas desaparecieron; el cluster de economía ahora es honesto
        (\emph{usgs, old scrap} = reciclaje/commodities).
  \item Emergió un \textbf{4º cluster} que estaba enterrado: literatura \textbf{en francés}
        sobre recuperación de Nb-Ta-Sn --- probablemente la escuela francesa que estudia
        Beauvoir. Hallazgo valioso... y lección nueva: \emph{el idioma también agrupa}
        (SPECTER junta los papers franceses). Refinamiento futuro: ¿traducir abstracts o
        aceptar ese cluster como subtema válido?
  \item El ruido subió (9 $\rightarrow$ 21): al re-acomodarse el mapa, más papers quedaron
        entre nubes. Para eso está tu slider de granularidad.
\end{itemize}

\begin{ojo}
Así se ve el ciclo real del data science: \textbf{construir $\rightarrow$ mirar $\rightarrow$
detectar algo raro $\rightarrow$ diagnosticar $\rightarrow$ arreglar $\rightarrow$ descubrir
algo nuevo}. Tu ojo experto detectó el problema; los embeddings dieron la solución; y el
arreglo destapó un hallazgo (el cluster francés) que nadie buscaba. Cada iteración enseña.
\end{ojo}

\begin{resumen}
Relevancia $=$ coseno(paper, TOPIC); umbral 0.55 elegido por el hueco natural en los datos;
filtrado antes de UMAP. Resultado: mapa limpio, etiquetas honestas y un cluster francés
sorpresa sobre Beauvoir. El sistema se limpia con sus propias herramientas.
\end{resumen}

\end{document}
